% ==============================================================
% Kapitel 1: Ziele und Aufgabenstellung
% ==============================================================
\section{Einleitung}
\subsection{Ziele}

Das Ziel dieser Diplomarbeit ist es, ein wartbares, erweiterbares und einfach 
aufsetzbares Zeitmesssystem für die \textit{Zero Emission Challenge} zu entwickeln, 
welches den gesamten Ablauf von der Zeiterfassung über die Punkteauswertung bis zur Ergebnisanzeige umfasst.
Das System soll zusätzlich auch die Verwaltung von Teams und Benutzern ermöglichen und 
eine rollenbasierte Zugriffskontrolle implementieren.

Das entwickelte System soll:
\begin{itemize}
    \item Renndaten aller Bewerbe zuverlässig erfassen, speichern und auswerten
      \item Versuchsdaten per API in das System speichern zu können
    \item Eine webbasierte Plattform zur Verwaltung von Teams, Benutzern und 
          Ergebnissen bereitstellen
    \item Sicherheit durch rollenbasierte Zugriffskontrolle und eine zentrale 
          Authentifizierung über Keycloak gewährleisten
    \item Mittels Docker einfach und plattformunabhängig betreibbar sein
      \item Softwarequalität durch automatisierte Tests wie Pytest und Playwright sicherstellen
\end{itemize}
\newpage

\subsection{Individuelle Aufgabenstellung}

Um diesen Zielen gerecht zu werden, wurden folgende Aufgaben definiert, die im Rahmen dieser Diplomarbeit umgesetzt werden sollen:

\begin{itemize}
      \item Entwicklung eines in Microservices aufgeteilten RESTful-API-Backends zur Zeiterfassung, Punkteauswertung sowie Team- und Benutzerverwaltung
      \item Entwicklung einer webbasierten Anwendung mit unterschiedlichen Funktionen für die Benutzergruppen Administrator, Teamleiter und Zuschauer
      \item Implementierung eines dedizierten Authentifizierungsdienstes
      \item Rollenbasierte Zugriffskontrolle auf die Website sowie die API-Endpunkte
      \item Modellierung und Anbindung einer relationalen PostgreSQL-Datenbank zur Datenspeicherung
      \item Containerisierung aller Systemkomponenten mittels Docker
      \item Implementierung automatisierter Testfälle mit Pytest und Playwright
\end{itemize}